% L04_full.tex -- Lecture 4: Fintech Security and Regulation (Full Variant)
% Frames: 31 | Charts: 14 \includegraphics | Arc: 10-role
% Generated for: Financial Technology (FinTech) -- MSc Course, Spring 2026
\input{../../_shared/preamble.tex}

\subtitle{Navigating Compliance in the Digital Finance Era}

\begin{document}

% =============================================================================
% === WHY === (Frames 1--4: Title, Opening Cartoon, Learning Objectives, Bridge)
% =============================================================================

% --- Frame 1: Title Page ---
\begin{frame}{Title Page}
  \titlepage
  \bottomnote{Lecture~4 of 7 --- Financial Technology (FinTech), Spring 2026.}
\end{frame}

% --- Frame 2: Opening Cartoon ---
\begin{frame}{``The Compliance Department's New Hire''}
  \begin{center}
    \includegraphics[width=0.85\textwidth]{11_opening_cartoon/cartoon.pdf}
  \end{center}
  \vspace{-0.5em}
  \begin{center}
    \textit{``It read all 27,000 pages of financial regulation in 3.2
    seconds. We're still not sure if it understood the footnotes.''}
  \end{center}
  \bottomnote{The EU financial services acquis exceeds 27,000 pages. The US Code of Federal Regulations, Title 12 (Banks and Banking), runs to over 9,000 pages --- and grows annually.}
\end{frame}

% --- Frame 3: Learning Objectives ---
\begin{frame}{Learning Objectives}
  By the end of this lecture you will be able to:
  \begin{enumerate}
    \item \textbf{Describe} the major regulatory approaches to fintech
          (innovation-friendly vs.\ precautionary) and the institutional
          architecture of key jurisdictions. \hfill\textit{[Understand]}
    \item \textbf{Explain} anti-money laundering (AML) and know-your-customer
          (KYC) processes, including the three stages of money laundering and
          digital identity verification. \hfill\textit{[Understand]}
    \item \textbf{Apply} a cost-analysis framework to the compliance burden
          facing fintechs and assess how RegTech reduces marginal compliance
          costs. \hfill\textit{[Apply]}
    \item \textbf{Analyze} regulatory arbitrage strategies across jurisdictions
          and evaluate the resulting fragmentation in digital asset
          regulation. \hfill\textit{[Analyze]}
    \item \textbf{Evaluate} the design trade-offs of regulatory sandboxes,
          weighing innovation acceleration against consumer protection
          risks. \hfill\textit{[Evaluate]}
  \end{enumerate}
  \vspace{0.5em}
  \textcolor{mlpurple}{\textbf{Bloom's levels covered:}} Understand, Apply, Analyze, Evaluate
  \bottomnote{These objectives map directly to quiz and exercise assessments.}
\end{frame}

% --- Frame 4: Bridge from Lecture 3 ---
\begin{frame}{Bridge from Lecture~3}
  \begin{columns}[T]
    \begin{column}{0.55\textwidth}
      In Lecture~3 we explored \textbf{payment infrastructure}: the four-party
      model, real-time rails, cross-border complexity, and interchange
      economics. We touched on payment-specific regulation --- PSD2 in the EU,
      the Durbin Amendment in the US, global interchange caps.

      \vspace{0.5em}
      Now we \textcolor{mlpurple}{\textbf{zoom out}} from payment regulation to
      the full regulatory landscape governing fintech:
      \begin{itemize}
        \item \textbf{From interchange caps to AML:} L03 showed how regulation
              shapes cost. L04 shows how regulation shapes \textit{access} ---
              who can operate, who can transact, and who is surveilled.
        \item \textbf{From PSD2 to MiCA:} Open banking was sector-specific.
              Comprehensive crypto regulation (MiCA) is economy-wide.
        \item \textbf{From settlement risk to systemic risk:} L03's settlement
              delays are micro-risks. L04's regulatory failures are macro-risks.
      \end{itemize}

      \vspace{0.3em}
      L04 shifts from \textcolor{mlpurple}{payment mechanics} to
      \textcolor{mlpurple}{the regulatory architecture that governs all of fintech}.
    \end{column}
    \begin{column}{0.42\textwidth}
      \includegraphics[width=\textwidth]{01_regulatory_approaches/chart.pdf}
    \end{column}
  \end{columns}
  \bottomnote{L03 gave you the payment infrastructure. L04 shows you the rules, enforcers, and technologies that determine who is allowed to build on that infrastructure.}
\end{frame}

% =============================================================================
% === FEEL === (Frame 5: Personal Connection)
% =============================================================================

% --- Frame 5: The Surveillance Paradox ---
\begin{frame}{The Surveillance Paradox}
  Think about what your bank knows about you. Every purchase, every transfer,
  every paycheck. Your rent, your subscriptions, your midnight food deliveries.
  \textcolor{mlpurple}{Your bank knows more about your daily life than your
  doctor, your employer, or your closest friend.}

  \vspace{0.5em}
  Is that \textbf{protection} or \textbf{surveillance}?

  \vspace{0.3em}
  \begin{itemize}
    \item \textcolor{mlpurple}{\textbf{The protection argument:}} Transaction
          monitoring catches fraud, money laundering, and terrorist financing.
          Without surveillance, criminals exploit the system, and ordinary
          consumers bear the cost through higher fees and systemic risk.
    \item \textcolor{mlpurple}{\textbf{The surveillance argument:}} Financial
          data reveals political donations, religious affiliations, health
          conditions, and personal relationships. Governments have historically
          weaponized financial surveillance --- from Nixon's ``enemies list'' to
          China's social credit system.
    \item \textcolor{mlpurple}{\textbf{The paradox:}} The same data that
          protects you from fraud \textit{can be used to control you}. The
          difference is governance, not technology.
  \end{itemize}

  \vspace{0.3em}
  \begin{exampleblock}{Quick Exercise}
    Open your banking app. Scroll through three months of transactions. What
    could a government learn about you? What could an employer? A landlord?
  \end{exampleblock}
  \bottomnote{Operation Choke Point (US, 2013--2017) used bank compliance pressure to de-bank legal businesses the government disfavored. Financial surveillance is never purely neutral.}
\end{frame}

% =============================================================================
% === WHAT === (Frames 6--9: Foundational Concepts -- Regulatory Philosophy)
% =============================================================================

% --- Frame 6: Innovation-Friendly vs. Precautionary ---
\begin{frame}{Innovation-Friendly vs.\ Precautionary Regulation}
  Two philosophical poles define the global approach to fintech regulation:

  \vspace{0.5em}
  \begin{columns}[T]
    \begin{column}{0.50\textwidth}
      \textcolor{mlpurple}{\textbf{Innovation-Friendly (Sandbox-First):}}
      \begin{itemize}
        \item \textbf{Philosophy:} Regulate outcomes, not activities. Let
              innovators experiment under supervision. Intervene only when
              harm materializes.
        \item \textbf{Jurisdictions:} UK (FCA), Singapore (MAS), Australia
              (ASIC), Abu Dhabi (ADGM).
        \item \textbf{Tools:} Regulatory sandboxes, innovation hubs, no-action
              letters, phased licensing.
        \item \textbf{Risk:} Under-protection of early consumers. Sandbox
              graduates may outgrow their regulatory training wheels.
      \end{itemize}
    \end{column}
    \begin{column}{0.47\textwidth}
      \textcolor{mlpurple}{\textbf{Precautionary (Rule-First):}}
      \begin{itemize}
        \item \textbf{Philosophy:} Regulate activities comprehensively before
              permitting them. Prevent harm before it occurs. Err on the side
              of restriction.
        \item \textbf{Jurisdictions:} EU (MiCA/DORA), China (PBOC), India
              (RBI for crypto), New York (BitLicense).
        \item \textbf{Tools:} Ex-ante licensing, activity-based regulation,
              product bans, capital requirements.
        \item \textbf{Risk:} Innovation flight. Startups relocate to
              friendlier jurisdictions. Regulatory moats protect incumbents.
      \end{itemize}
    \end{column}
  \end{columns}

  \vspace{0.3em}
  \begin{block}{The Uncomfortable Truth}
    Neither approach has proven definitively superior. The UK sandbox produced
    Revolut; the EU's precautionary approach prevented a TerraUST-scale
    consumer loss. \textcolor{mlpurple}{Context determines which philosophy
    is appropriate.}
  \end{block}
  \bottomnote{Zetzsche et al.\ (2017) classify fintech regulation into ``same risk, same rules'' (EU) vs.\ ``innovation-facilitation'' (UK/Singapore) paradigms.}
\end{frame}

% --- Frame 7: Regulatory Objectives ---
\begin{frame}{Regulatory Objectives}
  \begin{columns}[T]
    \begin{column}{0.55\textwidth}
      Every financial regulator pursues four objectives --- and they
      \textbf{inevitably conflict}:

      \vspace{0.4em}
      \begin{enumerate}
        \item \textcolor{mlpurple}{\textbf{Financial stability:}} Prevent
              systemic failures that cascade through the economy. Requires
              capital buffers, stress testing, and activity restrictions ---
              all of which raise costs and slow innovation.
        \item \textcolor{mlpurple}{\textbf{Consumer protection:}} Ensure fair
              treatment, transparency, and recourse. Requires disclosure
              mandates, suitability rules, and complaint mechanisms ---
              which add compliance burden.
        \item \textcolor{mlpurple}{\textbf{Competition:}} Prevent monopolies
              and ensure market access for new entrants. Requires open APIs,
              interoperability mandates, and anti-trust enforcement ---
              which may conflict with stability goals.
        \item \textcolor{mlpurple}{\textbf{Innovation:}} Enable new products
              and business models. Requires regulatory flexibility and
              experimentation space --- which may conflict with consumer
              protection.
      \end{enumerate}
    \end{column}
    \begin{column}{0.42\textwidth}
      \includegraphics[width=\textwidth]{01_regulatory_approaches/chart.pdf}

      \vspace{0.5em}
      \begin{alertblock}{The Conflict}
        Maximizing consumer protection (strict licensing, disclosure) raises
        barriers to entry, reducing competition. Maximizing innovation
        (light-touch rules) increases systemic risk.
        \textcolor{mlred}{No regulator can maximize all four simultaneously.}
      \end{alertblock}
    \end{column}
  \end{columns}
  \bottomnote{The Basel Committee's core principles emphasize stability; the IOSCO principles emphasize investor protection; the FSB emphasizes both. No framework weights innovation equally.}
\end{frame}

% --- Frame 8: The Regulatory Trilemma ---
\begin{frame}{The Regulatory Trilemma}
  Like Mundell's impossible trinity in macroeconomics, fintech regulation faces
  its own trilemma. A jurisdiction can optimize for \textbf{two of three}
  objectives, but not all three simultaneously:

  \vspace{0.5em}
  \begin{center}
    \begin{tikzpicture}[scale=1.3]
      % Triangle vertices
      \coordinate (A) at (0, 3.2);     % Innovation Speed
      \coordinate (B) at (-2.8, -0.5); % Consumer Safety
      \coordinate (C) at (2.8, -0.5);  % Financial Stability

      % Triangle edges
      \draw[thick, MLTEAL] (A) -- (B) -- (C) -- cycle;

      % Vertex labels
      \node[above=4pt, font=\bfseries\small, MLTEAL] at (A)
        {Innovation Speed};
      \node[below left=4pt, font=\bfseries\small, MLTEAL] at (B)
        {Consumer Safety};
      \node[below right=4pt, font=\bfseries\small, MLTEAL] at (C)
        {Financial Stability};

      % Edge labels (trade-offs)
      \node[left=6pt, font=\scriptsize, rotate=58, MLPURPLE] at
        ($(A)!0.5!(B)$) {Sandboxes (UK, SG)};
      \node[right=6pt, font=\scriptsize, rotate=-58, MLPURPLE] at
        ($(A)!0.5!(C)$) {Charter flexibility (US fintech)};
      \node[below=6pt, font=\scriptsize, MLPURPLE] at
        ($(B)!0.5!(C)$) {Precautionary rules (EU MiCA)};

      % Center note
      \node[font=\scriptsize\itshape, text width=3.2cm, align=center] at
        (0, 0.8) {Pick any two ---\\the third suffers};
    \end{tikzpicture}
  \end{center}

  \vspace{0.3em}
  \begin{block}{Implication}
    \textcolor{mlpurple}{Regulatory design is not a technical problem --- it
    is a political choice about which values a society prioritizes.} The
    ``right'' answer depends on the jurisdiction's risk tolerance, market
    maturity, and institutional capacity.
  \end{block}
  \bottomnote{Adapted from Brummer and Yadav (2019). The original monetary policy trilemma (Mundell, 1963) demonstrates the same logic: you cannot simultaneously control exchange rates, monetary policy, and capital flows.}
\end{frame}

% --- Frame 9: Global Regulatory Spectrum ---
\begin{frame}{Global Regulatory Spectrum}
  \begin{columns}[T]
    \begin{column}{0.50\textwidth}
      \includegraphics[width=\textwidth]{06_global_regulatory_comparison/chart.pdf}
    \end{column}
    \begin{column}{0.47\textwidth}
      Where major jurisdictions sit on the innovation--caution spectrum:

      \vspace{0.4em}
      \begin{itemize}
        \item \textbf{Most innovation-friendly:} Singapore (MAS), UAE (ADGM),
              UK (FCA). Sandbox-first, outcome-based, iterative licensing.
        \item \textbf{Balanced approach:} Australia (ASIC), Hong Kong (SFC/HKMA),
              Switzerland (FINMA). Clear rules with innovation windows.
        \item \textbf{Precautionary:} EU (MiCA/DORA), Japan (FSA). Comprehensive
              ex-ante regulation. High compliance burden but strong consumer
              protection.
        \item \textbf{Most restrictive:} China (outright crypto ban since 2021),
              India (RBI resistance to crypto, high taxation), New York
              (BitLicense regime).
        \item \textbf{Fragmented:} The United States. Federal agencies disagree
              on jurisdiction. 50 states add 50 licensing regimes. No unified
              fintech charter exists.
      \end{itemize}
    \end{column}
  \end{columns}
  \bottomnote{The Global Financial Innovation Network (GFIN) connects 80+ regulators to share sandbox learnings. Despite this, regulatory convergence remains limited.}
\end{frame}

% =============================================================================
% === CASE === (Frames 10--13: AML/KYC Deep Dive)
% =============================================================================

% --- Frame 10: Three Stages of Money Laundering ---
\begin{frame}{Three Stages of Money Laundering}
  \begin{columns}[T]
    \begin{column}{0.50\textwidth}
      \includegraphics[width=\textwidth]{04_money_laundering_stages/chart.pdf}
    \end{column}
    \begin{column}{0.47\textwidth}
      Money laundering converts illicit proceeds into apparently legitimate
      funds through three stages:

      \vspace{0.4em}
      \begin{enumerate}
        \item \textcolor{mlpurple}{\textbf{Placement}} --- Introducing dirty
              money into the financial system. Methods: structuring (smurfing)
              deposits below reporting thresholds, cash-intensive businesses
              (car washes, restaurants), currency exchange. \textit{Detection:}
              relatively easy --- cash transactions trigger automatic reports.
        \item \textcolor{mlpurple}{\textbf{Layering}} --- Obscuring the money
              trail through complex transactions. Methods: shell companies,
              offshore transfers, trade-based laundering, cryptocurrency mixing.
              \textit{Detection:} difficult --- requires cross-institution and
              cross-border data correlation.
        \item \textcolor{mlpurple}{\textbf{Integration}} --- Re-entering the
              legitimate economy. Methods: real estate purchases, luxury goods,
              legitimate business investment, loan-back schemes.
              \textit{Detection:} extremely difficult --- funds appear clean.
      \end{enumerate}
    \end{column}
  \end{columns}
  \bottomnote{The UN Office on Drugs and Crime estimates that 2--5\% of global GDP (USD~800B--2T annually) is laundered. Less than 1\% of illicit financial flows are seized.}
\end{frame}

% --- Frame 11: Know Your Customer (KYC) ---
\begin{frame}{Know Your Customer (KYC)}
  \begin{columns}[T]
    \begin{column}{0.50\textwidth}
      \includegraphics[width=\textwidth]{03_kyc_process_flow/chart.pdf}
    \end{column}
    \begin{column}{0.47\textwidth}
      KYC is the foundation of AML compliance --- the process by which
      financial institutions verify and monitor customer identity:

      \vspace{0.4em}
      \begin{itemize}
        \item \textbf{Customer Identification Program (CIP):} Collect and
              verify name, date of birth, address, and government-issued ID.
              Required at account opening.
        \item \textbf{Customer Due Diligence (CDD):} Assess the customer's
              risk profile based on occupation, source of funds, transaction
              patterns, and geographic exposure.
        \item \textbf{Enhanced Due Diligence (EDD):} Triggered for
              high-risk customers: politically exposed persons (PEPs),
              high-net-worth individuals, customers in sanctioned jurisdictions.
        \item \textbf{Ongoing Monitoring:} Continuous screening against
              sanctions lists, adverse media, and transaction anomalies.
              KYC is not a one-time event --- it is a lifecycle obligation.
      \end{itemize}

      \vspace{0.3em}
      \textcolor{mlpurple}{Average KYC onboarding cost: USD~25--50 per retail
      customer; USD~2,000--25,000 per corporate customer.}
    \end{column}
  \end{columns}
  \bottomnote{FATF Recommendation 10 requires all member countries to implement CDD. The 40 FATF Recommendations form the global AML/KYC baseline adopted by 200+ jurisdictions.}
\end{frame}

% --- Frame 12: Digital Identity Verification ---
\begin{frame}{Digital Identity Verification}
  Traditional KYC required in-person document inspection. \textbf{eKYC}
  (electronic KYC) replaces this with technology-driven verification:

  \vspace{0.5em}
  \begin{columns}[T]
    \begin{column}{0.50\textwidth}
      \textcolor{mlpurple}{\textbf{Core eKYC Technologies:}}
      \begin{itemize}
        \item \textbf{Document AI:} Optical character recognition (OCR) and
              machine learning extract data from passports, driver's licenses,
              and utility bills. Detects tampering, forgery, and expired
              documents automatically.
        \item \textbf{Biometric matching:} Facial recognition compares a live
              selfie to the document photo. Accuracy exceeds 99.5\% for
              leading vendors (NIST FRVT benchmark).
        \item \textbf{Liveness detection:} Prevents spoofing via printed photos,
              video replays, or deepfakes. Active (blink, turn head) and
              passive (texture analysis) methods.
        \item \textbf{Database verification:} Cross-references against
              government registries, credit bureaus, and sanctions databases
              in real time.
      \end{itemize}
    \end{column}
    \begin{column}{0.47\textwidth}
      \textcolor{mlpurple}{\textbf{Impact on Fintech Onboarding:}}

      \vspace{0.5em}
      \begin{itemize}
        \item Traditional bank onboarding: 2--4 weeks, branch visit required
        \item eKYC onboarding: 3--10 minutes, fully remote
        \item Cost reduction: 60--80\% vs.\ manual KYC processes
        \item False rejection rates: 3--8\% (bias concerns for darker skin
              tones and older adults)
      \end{itemize}

      \vspace{0.3em}
      \begin{alertblock}{The Inclusion Tension}
        eKYC assumes smartphone access, government-issued ID, and digital
        literacy. The \textcolor{mlred}{1.4~billion people} globally who lack
        official identification (World Bank ID4D) are excluded by the very
        technology designed to expand access.
      \end{alertblock}
    \end{column}
  \end{columns}
  \bottomnote{India's Aadhaar system (1.3B+ enrolled) demonstrates both the promise and peril of digital identity: it enabled eKYC at scale but raised profound privacy concerns.}
\end{frame}

% --- Frame 13: Transaction Monitoring and SAR Filing ---
\begin{frame}{Transaction Monitoring and SAR Filing}
  \begin{columns}[T]
    \begin{column}{0.50\textwidth}
      \includegraphics[width=\textwidth]{02_aml_compliance_flow/chart.pdf}
    \end{column}
    \begin{column}{0.47\textwidth}
      The end-to-end AML compliance pipeline runs continuously:

      \vspace{0.4em}
      \begin{enumerate}
        \item \textcolor{mlpurple}{\textbf{Data ingestion:}} Transaction data,
              customer profiles, external watchlists, and adverse media are
              aggregated in real time.
        \item \textcolor{mlpurple}{\textbf{Rule-based screening:}} Threshold
              alerts (e.g., transactions above USD~10,000), velocity checks,
              geographic risk flags, and sanctions list matching.
        \item \textcolor{mlpurple}{\textbf{ML-based anomaly detection:}}
              Behavioral profiling identifies deviations from normal patterns.
              Reduces false positives by 40--70\% vs.\ rules alone.
        \item \textcolor{mlpurple}{\textbf{Alert investigation:}} Compliance
              analysts review flagged transactions, gather supporting evidence,
              and make disposition decisions.
        \item \textcolor{mlpurple}{\textbf{SAR filing:}} Suspicious Activity
              Reports are submitted to the Financial Intelligence Unit (FIU).
              In the US, FinCEN receives over 4~million SARs annually.
      \end{enumerate}

      \vspace{0.3em}
      \textcolor{mlpurple}{\textbf{The false positive problem:}} 95--98\% of
      alerts are false positives. Each costs USD~20--50 to investigate manually.
    \end{column}
  \end{columns}
  \bottomnote{FinCEN received 4.6M SARs in 2023. The UK FIU received 900K+. The sheer volume overwhelms investigative capacity, reducing the signal-to-noise ratio.}
\end{frame}

% =============================================================================
% === HOW === (Frames 14--17: US Regulatory Architecture)
% =============================================================================

% --- Frame 14: US Regulation -- Federal Layer ---
\begin{frame}{US Regulation: Federal Layer}
  \begin{columns}[T]
    \begin{column}{0.50\textwidth}
      \includegraphics[width=\textwidth]{05_us_regulatory_patchwork/chart.pdf}
    \end{column}
    \begin{column}{0.47\textwidth}
      The US has \textbf{no single financial regulator}. Overlapping federal
      agencies create a complex jurisdictional patchwork:

      \vspace{0.4em}
      \begin{itemize}
        \item \textbf{OCC} (Office of the Comptroller of the Currency):
              Charters and supervises national banks. Proposed fintech charter
              (2018) stalled by state opposition.
        \item \textbf{SEC} (Securities and Exchange Commission): Regulates
              securities. Claims jurisdiction over most crypto tokens.
              Enforcement-first approach under Chair Gensler (2021--2025).
        \item \textbf{CFTC} (Commodity Futures Trading Commission): Regulates
              commodity derivatives. Claims Bitcoin and Ether as commodities.
        \item \textbf{CFPB} (Consumer Financial Protection Bureau): Supervises
              consumer financial products. Active in BNPL, earned wage access,
              and data privacy.
        \item \textbf{FinCEN} (Financial Crimes Enforcement Network): AML/BSA
              compliance. Issues rules on money transmission, virtual currency
              reporting, and beneficial ownership.
      \end{itemize}
    \end{column}
  \end{columns}
  \bottomnote{The US dual banking system dates to the National Bank Act (1863). Fintechs must navigate both federal and state regulators --- a burden that favors well-funded incumbents.}
\end{frame}

% --- Frame 15: US Regulation -- State Layer ---
\begin{frame}{US Regulation: State Layer}
  Below the federal layer, fintechs face a \textbf{50-state patchwork} of
  licensing requirements:

  \vspace{0.5em}
  \begin{columns}[T]
    \begin{column}{0.50\textwidth}
      \textcolor{mlpurple}{\textbf{Money Transmitter Licenses (MTLs):}}
      \begin{itemize}
        \item Required in 49 states + DC for any entity that transmits money
              (Montana is the sole exception).
        \item Each state has its own application process, fees (USD~500 to
              USD~100,000+), surety bond requirements, and examination
              schedules.
        \item Full 50-state licensing costs USD~2--5 million and takes
              12--24 months. This is a \textbf{de facto barrier to entry}
              for early-stage fintechs.
        \item NMLS (Nationwide Multistate Licensing System) provides a
              common portal but does not harmonize requirements.
      \end{itemize}
    \end{column}
    \begin{column}{0.47\textwidth}
      \textcolor{mlpurple}{\textbf{The OCC Fintech Charter Debate:}}
      \begin{itemize}
        \item In 2018, the OCC proposed a special-purpose national bank
              charter for fintechs --- preempting state MTL requirements.
        \item \textbf{State opposition:} The Conference of State Bank
              Supervisors (CSBS) sued, arguing federal overreach. New York
              DFS filed separately.
        \item \textbf{Result:} The charter was effectively shelved. Fintechs
              continue to navigate the state-by-state regime.
      \end{itemize}

      \vspace{0.3em}
      \begin{alertblock}{Regulatory Arbitrage}
        Some fintechs partner with state-chartered ``sponsor banks'' to
        avoid direct licensing. This \textcolor{mlred}{rent-a-charter} model
        concentrates risk in a small number of sponsoring banks.
      \end{alertblock}
    \end{column}
  \end{columns}
  \bottomnote{New York's BitLicense (2015) is the strictest state-level crypto regime. Only 35 BitLicenses have been issued in 10 years --- many firms simply do not serve New York.}
\end{frame}

% --- Frame 16: SEC vs.\ CFTC -- Digital Asset Divide ---
\begin{frame}{SEC vs.\ CFTC: Digital Asset Divide}
  The fundamental question in US digital asset regulation: \textbf{is a token
  a security or a commodity?} The answer determines which agency has
  jurisdiction --- and which rules apply.

  \vspace{0.5em}
  \begin{columns}[T]
    \begin{column}{0.50\textwidth}
      \textcolor{mlpurple}{\textbf{SEC: The Howey Test}}
      \begin{itemize}
        \item Under \textit{SEC v.\ W.J.\ Howey Co.} (1946), an ``investment
              contract'' exists when there is (1) an investment of money,
              (2) in a common enterprise, (3) with expectation of profits,
              (4) derived from the efforts of others.
        \item The SEC argues most crypto tokens satisfy Howey and are therefore
              unregistered securities.
        \item \textbf{Key enforcement:} SEC sued Ripple (XRP), Coinbase,
              Binance, and Kraken. Mixed judicial outcomes.
        \item Consequence: tokens classified as securities face registration
              requirements that are practically impossible for decentralized
              projects.
      \end{itemize}
    \end{column}
    \begin{column}{0.47\textwidth}
      \textcolor{mlpurple}{\textbf{CFTC: Commodity Classification}}
      \begin{itemize}
        \item The CFTC has classified Bitcoin and Ether as commodities,
              giving it jurisdiction over derivatives and spot-market fraud.
        \item Commodity classification imposes lighter regulatory burden:
              no registration for spot trading, no prospectus requirement.
        \item \textbf{The gap:} Neither agency clearly regulates spot crypto
              trading that does not involve securities or derivatives.
      \end{itemize}

      \vspace{0.3em}
      \begin{block}{Legislative Proposals}
        FIT21 (Financial Innovation and Technology for the 21st Century Act)
        proposed a clear SEC/CFTC dividing line based on decentralization.
        \textcolor{mlpurple}{As of 2026, comprehensive legislation remains
        pending.}
      \end{block}
    \end{column}
  \end{columns}
  \bottomnote{Judge Torres in SEC v.\ Ripple (2023) ruled that programmatic XRP sales on exchanges were not securities --- the first major judicial pushback on SEC's expansive claims.}
\end{frame}

% --- Frame 17: CFPB and Consumer Protection ---
\begin{frame}{CFPB and Consumer Protection}
  The Consumer Financial Protection Bureau (est.\ 2011, Dodd-Frank Act) is the
  primary US federal regulator for consumer financial products:

  \vspace{0.5em}
  \begin{columns}[T]
    \begin{column}{0.50\textwidth}
      \textcolor{mlpurple}{\textbf{Fintech Enforcement Areas:}}
      \begin{itemize}
        \item \textbf{Buy Now, Pay Later (BNPL):} In 2024, the CFPB issued an
              interpretive rule classifying BNPL providers as ``card issuers''
              under Regulation Z, requiring dispute rights and billing
              disclosures.
        \item \textbf{Earned Wage Access (EWA):} Providers like Earnin and
              DailyPay offer early paycheck access. The CFPB is evaluating
              whether ``tips'' and ``express fees'' constitute hidden interest.
        \item \textbf{Data privacy (Section 1033):} The CFPB's Personal
              Financial Data Rights rule (2024) mandates open banking-style
              data portability --- the US equivalent of PSD2's access provisions.
      \end{itemize}
    \end{column}
    \begin{column}{0.47\textwidth}
      \textcolor{mlpurple}{\textbf{The Political Dimension:}}
      \begin{itemize}
        \item The CFPB's authority has been contested since inception. The
              Supreme Court upheld its funding mechanism in
              \textit{CFPB v.\ Community Financial Services} (2024).
        \item Enforcement intensity varies with administration: aggressive
              under Democratic leadership, curtailed under Republican
              leadership.
        \item \textbf{Fintech implication:} Regulatory uncertainty forces
              fintechs to design for the \textit{most restrictive} possible
              interpretation, increasing compliance costs.
      \end{itemize}

      \vspace{0.3em}
      \begin{alertblock}{The Pattern}
        US consumer protection regulation is \textcolor{mlred}{enforcement-driven
        rather than legislation-driven}: rules emerge from lawsuits, consent
        orders, and interpretive guidance --- not from statute.
      \end{alertblock}
    \end{column}
  \end{columns}
  \bottomnote{The CFPB has returned over USD~19.6B to consumers since 2011. Its enforcement actions increasingly target fintech firms alongside traditional banks.}
\end{frame}

% =============================================================================
% === RISK === (Frames 18--20: Global Regulatory Regimes)
% =============================================================================

% --- Frame 18: EU MiCA Regulation ---
\begin{frame}{EU MiCA Regulation}
  \begin{columns}[T]
    \begin{column}{0.55\textwidth}
      The Markets in Crypto-Assets Regulation (\textbf{MiCA}, effective
      June 2024 for stablecoins, December 2024 for all other provisions) is
      the world's first comprehensive crypto regulatory framework:

      \vspace{0.4em}
      \begin{itemize}
        \item \textbf{Token classification:} Three categories ---
              Asset-Referenced Tokens (ARTs, backed by multiple assets),
              E-Money Tokens (EMTs, backed by single fiat currency), and
              ``other'' crypto-assets (utility tokens, meme coins).
        \item \textbf{Licensing:} All Crypto-Asset Service Providers (CASPs)
              must obtain authorization from a national competent authority.
              Passportable across all 27 EU member states.
        \item \textbf{Stablecoin requirements:} EMTs must be issued by
              authorized credit institutions or e-money institutions. 1:1
              reserve backing in segregated accounts. Redemption at par value
              at any time.
        \item \textbf{Consumer protection:} Mandatory white papers, conflict
              of interest rules, custody requirements, market abuse
              prohibitions.
      \end{itemize}
    \end{column}
    \begin{column}{0.42\textwidth}
      \includegraphics[width=\textwidth]{06_global_regulatory_comparison/chart.pdf}

      \vspace{0.5em}
      \begin{block}{The Brussels Effect}
        MiCA is expected to become the \textcolor{mlpurple}{global de facto
        standard} for crypto regulation --- similar to GDPR for data privacy.
        Non-EU firms serving EU customers must comply, creating extraterritorial
        reach.
      \end{block}
    \end{column}
  \end{columns}
  \bottomnote{MiCA explicitly excludes DeFi, NFTs (unless fungible), and CBDCs from its scope. These gaps will likely require future legislative action.}
\end{frame}

% --- Frame 19: UK FCA Regulatory Sandbox ---
\begin{frame}{UK FCA Regulatory Sandbox}
  The UK Financial Conduct Authority launched the world's first regulatory
  sandbox in 2016, establishing the template that 50+ countries have since
  adopted:

  \vspace{0.5em}
  \begin{columns}[T]
    \begin{column}{0.50\textwidth}
      \textcolor{mlpurple}{\textbf{How the Sandbox Works:}}
      \begin{itemize}
        \item \textbf{Application:} Firms propose an innovative product or
              service, define the test parameters, and identify consumer
              risks.
        \item \textbf{Cohort selection:} The FCA admits 20--30 firms per
              cohort (biannual). Acceptance rate: approximately 40\%.
        \item \textbf{Testing:} Firms operate with real customers under
              relaxed regulatory requirements (e.g., reduced capital, modified
              disclosure) for 3--6 months.
        \item \textbf{Safeguards:} Customer consent required. Compensation
              schemes apply. FCA can terminate testing at any time.
        \item \textbf{Graduation:} Successful firms apply for full
              authorization. Approximately 80\% of sandbox firms remain
              operational after graduation.
      \end{itemize}
    \end{column}
    \begin{column}{0.47\textwidth}
      \textcolor{mlpurple}{\textbf{Outcomes and Criticism:}}
      \begin{itemize}
        \item Over 800 firms have participated across 12+ cohorts since 2016.
              Notable graduates: Revolut, Monzo, Bud Financial.
        \item \textbf{Criticism 1:} Selection bias --- the sandbox admits
              firms most likely to succeed regardless, inflating success
              metrics.
        \item \textbf{Criticism 2:} Limited scale --- sandbox customers
              are often early adopters, not representative of vulnerable
              populations.
        \item \textbf{Criticism 3:} Post-sandbox cliff --- graduating firms
              face full regulatory burden immediately, with no transition
              period.
      \end{itemize}

      \vspace{0.3em}
      \begin{block}{Evolution}
        The FCA introduced a ``Digital Sandbox'' (permanent, data-driven) in
        2023 and a ``Global Sandbox'' with international regulators.
      \end{block}
    \end{column}
  \end{columns}
  \bottomnote{Cornelli et al.\ (2020, BIS) found that sandbox participation increased subsequent funding by 15\% on average, though the causal mechanism is debated (signaling vs.\ support).}
\end{frame}

% --- Frame 20: Singapore MAS Framework ---
\begin{frame}{Singapore MAS Framework}
  The Monetary Authority of Singapore (MAS) has positioned the city-state as
  Asia's fintech hub through a multi-layered regulatory framework:

  \vspace{0.5em}
  \begin{columns}[T]
    \begin{column}{0.50\textwidth}
      \textcolor{mlpurple}{\textbf{Payment Services Act (PSA, 2020):}}
      \begin{itemize}
        \item Activity-based licensing: firms are licensed for specific
              payment services (e-money issuance, digital payment tokens,
              cross-border transfers, merchant acquisition).
        \item Three license tiers based on transaction volume:
              \textbf{Standard} (up to SGD~3M/month),
              \textbf{Major} (above SGD~3M/month), and
              \textbf{Money-Changing} (physical currency exchange only).
        \item AML/CFT obligations calibrated by tier --- lower burden for
              smaller operators.
      \end{itemize}

      \vspace{0.3em}
      \textcolor{mlpurple}{\textbf{Sandbox Express (2019):}}
      \begin{itemize}
        \item Pre-defined sandbox conditions for common business models.
              Eliminates case-by-case negotiation. Approval in 21 days.
        \item Covers insurance distribution, remittance, and recognized
              market operators.
      \end{itemize}
    \end{column}
    \begin{column}{0.47\textwidth}
      \textcolor{mlpurple}{\textbf{Digital Asset Regulation:}}
      \begin{itemize}
        \item MAS distinguishes between ``digital payment tokens'' (DPTs,
              regulated for AML only) and ``capital markets products''
              (securities, regulated comprehensively).
        \item Crypto exchanges must obtain a Major Payment Institution
              license. Retail crypto advertising is restricted.
        \item \textbf{Project Ubin/Orchid:} MAS's wholesale CBDC experiments
              exploring tokenized settlement.
      \end{itemize}

      \vspace{0.3em}
      \begin{block}{The Singapore Model}
        MAS combines \textcolor{mlpurple}{principle-based regulation}
        (technology-neutral, outcome-focused) with
        \textcolor{mlpurple}{active ecosystem building} (grants, innovation
        labs, international cooperation). The result: over 1,600 fintech
        firms and growing.
      \end{block}
    \end{column}
  \end{columns}
  \bottomnote{Singapore's FinTech Festival (organized by MAS) is the world's largest, attracting 60,000+ participants. MAS uses it as both a policy platform and an ecosystem accelerator.}
\end{frame}

% =============================================================================
% === WHERE === (Frames 21--23: RegTech and Compliance Technology)
% =============================================================================

% --- Frame 21: RegTech -- Automated Compliance ---
\begin{frame}{RegTech: Automated Compliance}
  \begin{columns}[T]
    \begin{column}{0.50\textwidth}
      \includegraphics[width=\textwidth]{07_regtech_stack_architecture/chart.pdf}
    \end{column}
    \begin{column}{0.47\textwidth}
      \textbf{RegTech} (Regulatory Technology) applies technology to regulatory
      compliance --- reducing cost, increasing accuracy, and enabling real-time
      reporting:

      \vspace{0.4em}
      \begin{itemize}
        \item \textbf{Data ingestion layer:} APIs connect to transaction
              systems, customer databases, external watchlists, and regulatory
              feeds.
        \item \textbf{Analytics engine:} Machine learning models for
              transaction monitoring, risk scoring, and anomaly detection.
              Graph analytics map entity relationships.
        \item \textbf{Regulatory intelligence:} NLP parses regulatory
              updates, maps them to internal controls, and generates
              gap analyses.
        \item \textbf{Reporting and audit:} Automated generation of
              regulatory filings (SARs, CTRs, FATCA/CRS reports). Immutable
              audit trails for examiner review.
      \end{itemize}

      \vspace{0.3em}
      \textcolor{mlpurple}{Global RegTech spending reached USD~12.3B in 2024
      and is projected to exceed USD~28B by 2028 (Juniper Research).}
    \end{column}
  \end{columns}
  \bottomnote{Key RegTech vendors: ComplyAdvantage (AML screening), Chainalysis (blockchain analytics), Suade Labs (regulatory reporting), Onfido (identity verification), Behavox (conduct surveillance).}
\end{frame}

% --- Frame 22: NLP for Regulatory Change ---
\begin{frame}{NLP for Regulatory Change}
  Financial regulations change constantly. In 2023 alone, Thomson Reuters
  Regulatory Intelligence tracked over \textbf{61,000 regulatory updates}
  globally --- an average of 167 per day.

  \vspace{0.5em}
  \begin{columns}[T]
    \begin{column}{0.50\textwidth}
      \textcolor{mlpurple}{\textbf{The Manual Problem:}}
      \begin{itemize}
        \item Compliance teams must monitor hundreds of regulatory sources
              across multiple jurisdictions.
        \item Each update must be assessed for relevance, mapped to internal
              policies, and translated into operational changes.
        \item Manual regulatory change management costs large banks
              USD~50--100M annually (McKinsey, 2023).
        \item Median time from regulatory publication to internal
              implementation: 6--12 months --- creating compliance gaps.
      \end{itemize}
    \end{column}
    \begin{column}{0.47\textwidth}
      \textcolor{mlpurple}{\textbf{NLP-Powered Solution:}}
      \begin{itemize}
        \item \textbf{Ingestion:} Automated scraping of regulatory gazettes,
              central bank publications, enforcement actions, and
              consultation papers.
        \item \textbf{Classification:} NLP models classify updates by topic,
              jurisdiction, entity type, and urgency.
        \item \textbf{Impact mapping:} Semantic matching links regulatory
              text to specific internal policies, controls, and business
              lines.
        \item \textbf{Alert generation:} Compliance officers receive
              prioritized alerts with summaries, affected controls, and
              recommended actions.
      \end{itemize}

      \vspace{0.3em}
      \begin{block}{LLM Impact}
        Large language models are transforming regulatory NLP from keyword
        matching to \textcolor{mlpurple}{contextual interpretation} ---
        understanding regulatory intent, not just regulatory text.
      \end{block}
    \end{column}
  \end{columns}
  \bottomnote{Ascent RegTech, Cube, and NICE Actimize are leading NLP-based regulatory change management platforms. Adoption remains concentrated in Tier~1 banks.}
\end{frame}

% --- Frame 23: The Compliance Cost Burden ---
\begin{frame}{The Compliance Cost Burden}
  \begin{columns}[T]
    \begin{column}{0.50\textwidth}
      \includegraphics[width=\textwidth]{09_compliance_cost_comparison/chart.pdf}
    \end{column}
    \begin{column}{0.47\textwidth}
      Compliance costs are \textbf{regressive}: they disproportionately burden
      smaller firms, creating a structural advantage for incumbents.

      \vspace{0.4em}
      \begin{itemize}
        \item \textbf{Large banks (over USD~50B assets):} Compliance costs
              average 4--6\% of total operating expenses. Dedicated teams of
              1,000+ compliance staff. Economies of scale apply.
        \item \textbf{Mid-size banks (USD~1--50B):} 6--10\% of operating
              expenses. Compliance is the fastest-growing cost center.
        \item \textbf{Small banks and fintechs (under USD~1B):} 10--15\% of
              operating expenses. A single regulatory change can require
              disproportionate system investment.
        \item \textbf{The math:} A startup spending USD~500K on compliance
              might spend 15\% of revenue. A megabank spending USD~1B might
              spend 2\% of revenue.
      \end{itemize}

      \vspace{0.3em}
      \begin{block}{RegTech as Equalizer}
        Cloud-based RegTech solutions offer SaaS compliance at
        \textcolor{mlpurple}{USD~10K--100K/year} --- potentially leveling
        the playing field for small firms.
      \end{block}
    \end{column}
  \end{columns}
  \bottomnote{The LexisNexis True Cost of Compliance report (2024) estimates global financial services compliance spending at USD~274B annually. AML compliance alone exceeds USD~50B.}
\end{frame}

% =============================================================================
% === IMPACT === (Frames 24--25: Sandboxes and SupTech)
% =============================================================================

% --- Frame 24: Regulatory Sandboxes Worldwide ---
\begin{frame}{Regulatory Sandboxes Worldwide}
  \begin{columns}[T]
    \begin{column}{0.50\textwidth}
      \includegraphics[width=\textwidth]{08_sandbox_adoption_timeline/chart.pdf}
    \end{column}
    \begin{column}{0.47\textwidth}
      Since the UK's 2016 launch, regulatory sandboxes have proliferated
      globally --- with significant design variations:

      \vspace{0.4em}
      \begin{itemize}
        \item \textbf{2016:} UK (FCA), Singapore (MAS), Australia (ASIC),
              Hong Kong (HKMA). First movers.
        \item \textbf{2017--2018:} Canada (CSA), Malaysia (BNM), Thailand
              (BOT), Bahrain (CBB), Abu Dhabi (ADGM). Rapid adoption.
        \item \textbf{2019--2021:} EU member states (individual), India (RBI
              --- limited scope), Japan (FSA), Brazil (BCB). Second wave.
        \item \textbf{2022--present:} Over 80 jurisdictions operate sandboxes.
              Cross-border sandboxes (GFIN) enable multi-jurisdiction testing.
      \end{itemize}

      \vspace{0.3em}
      \textbf{Design variations:}
      \begin{itemize}
        \item \textbf{Cohort vs.\ rolling:} UK uses cohorts; Singapore
              accepts applications continuously.
        \item \textbf{Scope:} Some restrict to specific sectors (payments,
              insurance); others are open-ended.
        \item \textbf{Duration:} 3 months (minimum) to 24 months (maximum).
      \end{itemize}
    \end{column}
  \end{columns}
  \bottomnote{The World Bank (2020) catalogued 73 sandbox initiatives globally. By 2024, this number exceeded 80, though quality and activity levels vary significantly.}
\end{frame}

% --- Frame 25: SupTech -- The Regulator's AI ---
\begin{frame}{SupTech: The Regulator's AI}
  \textbf{SupTech} (Supervisory Technology) is RegTech's mirror image: technology
  used \textit{by regulators} to supervise financial institutions more
  effectively.

  \vspace{0.5em}
  \begin{columns}[T]
    \begin{column}{0.50\textwidth}
      \textcolor{mlpurple}{\textbf{Core SupTech Applications:}}
      \begin{itemize}
        \item \textbf{Automated reporting:} Machine-readable regulatory
              reports (XBRL, JSON) replace manual submissions. The Bank of
              England's Digital Regulatory Reporting initiative processes
              filings in minutes, not weeks.
        \item \textbf{Market surveillance:} AI monitors trading patterns,
              detects market manipulation (spoofing, layering, wash trading),
              and identifies insider trading networks.
        \item \textbf{Stress testing:} Machine learning models simulate
              systemic scenarios (pandemic, market crash, cyber attack) faster
              and with more granularity than traditional models.
        \item \textbf{Network analysis:} Graph algorithms map interconnections
              between financial institutions, identifying contagion risks and
              too-connected-to-fail entities.
      \end{itemize}
    \end{column}
    \begin{column}{0.47\textwidth}
      \textcolor{mlpurple}{\textbf{Leading SupTech Programs:}}
      \begin{itemize}
        \item \textbf{MAS (Singapore):} Data analytics group uses NLP to
              analyze financial institution disclosures and detect anomalies.
        \item \textbf{De Nederlandsche Bank:} Deployed machine learning for
              AML supervision and pension fund risk assessment.
        \item \textbf{Bank of Italy:} Uses text mining on bank communications
              to detect early warning signals of distress.
        \item \textbf{SEC (US):} MIDAS system processes billions of trading
              records daily for market surveillance.
      \end{itemize}

      \vspace{0.3em}
      \begin{alertblock}{The Arms Race}
        As regulated firms adopt AI for compliance, regulators must adopt AI
        for supervision. \textcolor{mlred}{A technology gap between supervisor
        and supervised creates systemic blind spots.}
      \end{alertblock}
    \end{column}
  \end{columns}
  \bottomnote{The BIS (2019) found that 40\% of central banks and supervisors were actively exploring SupTech. By 2024, this figure exceeded 70\% (FSI survey).}
\end{frame}

% =============================================================================
% === SO WHAT === (Frames 26--27: Synthesis and Evaluation)
% =============================================================================

% --- Frame 26: Embedded Compliance and Harmonization ---
\begin{frame}{Embedded Compliance and Harmonization}
  \begin{columns}[T]
    \begin{column}{0.55\textwidth}
      The future of compliance is \textcolor{mlpurple}{\textbf{embedded,
      automated, and API-first}} --- compliance as infrastructure, not as
      overhead.

      \vspace{0.4em}
      \begin{itemize}
        \item \textbf{Compliance-as-a-Service (CaaS):} Startups like
              Alloy, Unit, and Sardine provide turnkey compliance
              infrastructure via API. A fintech can integrate KYC, AML
              monitoring, and regulatory reporting without building
              in-house capability.
        \item \textbf{Embedded KYC:} Identity verification triggered
              automatically during onboarding flows. The user uploads
              an ID; the API returns a risk score in seconds.
        \item \textbf{Programmable compliance:} Smart contracts encode
              regulatory rules directly. Transfer restrictions, investor
              accreditation checks, and reporting obligations execute
              automatically.
        \item \textbf{API-first regulation:} Regulators in Singapore and
              the UK are experimenting with machine-readable regulation ---
              rules published as code, not just as PDF documents.
      \end{itemize}
    \end{column}
    \begin{column}{0.42\textwidth}
      \includegraphics[width=\textwidth]{10_regtech_investment_trends/chart.pdf}

      \vspace{0.5em}
      \begin{block}{Regulatory Harmonization}
        Cross-border fintech requires cross-border rules. Emerging
        harmonization efforts:
        \begin{itemize}
          \item MiCA (EU-wide crypto framework)
          \item FATF travel rule (global crypto AML)
          \item Basel III/IV (global bank capital)
          \item GFIN (global sandbox cooperation)
        \end{itemize}
        \textcolor{mlpurple}{Full harmonization remains a distant goal ---
        but convergence is accelerating.}
      \end{block}
    \end{column}
  \end{columns}
  \bottomnote{RegTech venture funding exceeded USD~6B cumulatively through 2024. Compliance-as-a-Service is the fastest-growing segment within the broader RegTech market.}
\end{frame}

% --- Frame 27: Regulatory Evaluation Framework ---
\begin{frame}{Regulatory Evaluation Framework}
  \begin{columns}[T]
    \begin{column}{0.55\textwidth}
      Five questions to evaluate any fintech regulatory regime:

      \vspace{0.4em}
      \begin{enumerate}
        \item \textbf{Does it protect consumers effectively?}\\
              Are disclosure requirements meaningful? Is redress accessible?
              Are vulnerable populations considered? Or does ``protection''
              mean ``exclusion by compliance cost''?
        \item \textbf{Does it enable innovation?}\\
              Can a startup with a good idea get to market in a reasonable
              time and at reasonable cost? Or do licensing requirements
              create insurmountable barriers to entry?
        \item \textbf{Does it maintain financial stability?}\\
              Are capital requirements, stress testing, and resolution
              frameworks adequate? Or does light-touch regulation create
              systemic blind spots?
        \item \textbf{Does it scale?}\\
              Does the regime work for 10 firms and for 10,000 firms? For
              domestic players and for cross-border operators?
        \item \textbf{Does it adapt?}\\
              Can the regime respond to new technologies (AI, DeFi, quantum)
              without requiring new legislation for each? Or does it become
              obsolete between regulatory cycles?
      \end{enumerate}
    \end{column}
    \begin{column}{0.42\textwidth}
      \textbf{Applying the framework:}

      \vspace{0.4em}
      \begin{block}{UK FCA Sandbox}
        Consumer protection: moderate (safeguards but limited scale).
        Innovation: strong. Stability: untested at scale. Scalability:
        limited (cohort model). Adaptability: high.
      \end{block}

      \vspace{0.3em}
      \begin{block}{EU MiCA}
        Consumer protection: strong. Innovation: moderate (high compliance
        cost). Stability: strong (capital requirements). Scalability:
        strong (passporting). Adaptability: low (legislative process).
      \end{block}

      \vspace{0.3em}
      \begin{block}{US Patchwork}
        Consumer protection: variable by state. Innovation: weak (licensing
        burden). Stability: strong (bank-centric). Scalability: weak
        (50-state fragmentation). Adaptability: low.
      \end{block}
    \end{column}
  \end{columns}
  \bottomnote{This framework integrates L01's strategy lens, L03's payment regulation analysis, and L04's regulatory architecture into a single evaluation tool.}
\end{frame}

% =============================================================================
% === ACT === (Frame 28: Forward Look)
% =============================================================================

% --- Frame 28: What Comes Next ---
\begin{frame}{What Comes Next}
  \begin{itemize}
    \item \textbf{Next:} L05 (Personal Finance and Wealth Management) ---
          How technology is democratizing investment advice, automating
          portfolio management, and reshaping the advisor-client relationship.
          Robo-advisors, algorithmic trading, and the personalization of
          financial planning.
    \item \textbf{Before L05, reflect:} Pick a fintech product you use (a
          neobank, a payment app, a crypto exchange). Research which regulator
          supervises it and what license it holds. Were you surprised by the
          answer --- or by the lack of one?
    \item \textbf{Workshop preparation:} Review the regulatory evaluation
          framework (Frame~27). You will use it to compare two regulatory
          regimes in Workshop~E.
  \end{itemize}

  \vspace{0.8em}
  \begin{block}{Why Personal Finance Matters}
    Every regulation discussed today --- AML, KYC, sandbox, MiCA ---
    \textcolor{mlpurple}{exists to protect the individual consumer}. L05
    examines what that consumer actually does with their money: saving,
    investing, and planning for the future. The regulatory guardrails of L04
    enable the wealth management innovations of L05.
  \end{block}

  \vspace{0.5em}
  \begin{exampleblock}{Course Progress}
    \small
    L01: Foundations~\checkmark \(\bullet\) L02: Ecosystem~\checkmark
    \(\bullet\) L03: Payments~\checkmark \(\bullet\)
    \textbf{L04: Regulation}~\checkmark \(\bullet\)
    L05: Wealth Mgmt \(\bullet\) L06: Insurance \(\bullet\)
    L07: Technology
  \end{exampleblock}
  \bottomnote{L05 builds directly on L04: the regulatory frameworks governing investment advice (MiFID~II, Regulation Best Interest) are direct applications of the principles discussed today.}
\end{frame}

% =============================================================================
% === CLOSING === (Frames 29--31: Closing Cartoon, Key Takeaways, Summary)
% =============================================================================

% --- Frame 29: Closing Cartoon ---
\begin{frame}{``SupTech vs.\ Traditional Supervision''}
  \begin{center}
    \includegraphics[width=0.85\textwidth]{12_closing_cartoon/cartoon.pdf}
  \end{center}
  \vspace{-0.5em}
  \begin{center}
    \textit{``Our AI flagged 14 suspicious transactions while you were
    reviewing page 3 of your checklist.'' --- ``I found a typo on page 2.''}
  \end{center}
  \bottomnote{The BIS estimates that SupTech can reduce supervisory review times by 50--70\%. The challenge is not technology but institutional adoption and explainability requirements.}
\end{frame}

% --- Frame 30: Key Takeaways ---
\begin{frame}{Key Takeaways}
  \begin{enumerate}
    \item \textbf{Two regulatory philosophies:} Innovation-friendly (sandbox-first,
          outcome-based) vs.\ precautionary (rule-first, prevention-based).
          Neither has proven universally superior --- context determines the
          right approach.
    \item \textbf{AML/KYC is the compliance backbone:} Three laundering stages
          (placement, layering, integration) require layered defenses: CIP, CDD,
          EDD, and ongoing transaction monitoring. Less than 1\% of illicit flows
          are caught.
    \item \textbf{The US is uniquely fragmented:} Five federal agencies and 50
          state licensing regimes create the world's most complex fintech
          regulatory environment --- a de facto barrier to entry.
    \item \textbf{MiCA sets the global template:} The EU's comprehensive crypto
          framework will likely become the de facto global standard through the
          Brussels Effect.
    \item \textbf{RegTech reduces compliance costs by 60--80\%:} Automated KYC,
          ML-based transaction monitoring, and NLP-driven regulatory intelligence
          are transforming compliance from a cost center to a competitive advantage.
    \item \textbf{Sandboxes accelerate but do not replace regulation:} Over 80
          jurisdictions operate sandboxes, but graduation and scaling challenges
          remain.
    \item \textbf{SupTech closes the supervision gap:} Regulators are adopting
          the same AI tools they regulate, creating a technology arms race
          between supervisor and supervised.
  \end{enumerate}
  \bottomnote{Review question: Apply the five-question regulatory evaluation framework to the jurisdiction where you live. Where does it excel and where does it fall short?}
\end{frame}

% --- Frame 31: Summary and Key Vocabulary ---
\begin{frame}{Summary and Key Vocabulary}
  \textbf{Summary:} Financial regulation is the invisible architecture that
  determines who can build fintech products, who can use them, and who is
  protected when they fail. The global landscape ranges from innovation-friendly
  sandboxes (UK, Singapore) to comprehensive precautionary frameworks (EU MiCA)
  to fragmented patchworks (US). At the core of every regime lies AML/KYC
  compliance --- the processes that verify identity, monitor transactions, and
  report suspicious activity. RegTech is transforming compliance from a
  manual, costly burden into an automated, scalable capability, while SupTech
  gives regulators the tools to keep pace with the firms they supervise. The
  central lesson of L04 is that \textit{regulation is not an obstacle to
  fintech innovation --- it is the \textcolor{mlpurple}{framework within which
  sustainable innovation occurs}}. The jurisdictions that get the balance right
  will attract the next generation of financial technology.

  \vspace{0.5em}
  \textbf{Key Vocabulary:}
  \begin{multicols}{2}
    \begin{itemize}
      \item Regulatory Sandbox
      \item AML / KYC / CDD / EDD
      \item Money Laundering (Placement, Layering, Integration)
      \item Suspicious Activity Report (SAR)
      \item MiCA (Markets in Crypto-Assets)
      \item RegTech / SupTech
      \item Howey Test
      \item Regulatory Arbitrage
      \item Compliance-as-a-Service (CaaS)
      \item FATF Recommendations
    \end{itemize}
  \end{multicols}

  \vspace{0.3em}
  \textbf{Next lesson:} \textit{Lecture~5: Personal Finance and Wealth
  Management} --- robo-advisors, algorithmic portfolio management,
  financial planning technology, and the democratization of investment advice.
  \bottomnote{L05 connects L04's regulatory guardrails to the consumer-facing products they protect. MiFID~II, Reg~BI, and fiduciary standards are direct applications of today's principles.}
\end{frame}

\end{document}
